\begin{table}[htbp]\centering\footnotesize
\caption{Los siete municipios que sólo existen bajo la clave actualizada}\label{cua:zapnuevos}
\begin{tabular}{ll>{\raggedright\arraybackslash}p{2.4cm}>{\raggedright\arraybackslash}p{4.0cm}r}
\toprule
Clave & Municipio & Entidad & Se desprende de & AGEB \\
\midrule
02007 & San Felipe & Baja California & 02002 Mexicali & 16 \\
04013 & Dzitbalché & Campeche & 04001 Calkiní & 8 \\
12082 & Las Vigas & Guerrero & 12053 San Marcos & 8 \\
12085 & San Nicolás & Guerrero & 12023 Cuajinicuilapa & 4 \\
24059 & Villa de Pozos & San Luis Potosí & 24028 San Luis Potosí & 18 \\
25019 & Eldorado & Sinaloa & 25006 Culiacán & 25 \\
25020 & Juan José Ríos & Sinaloa & 25001 Ahome; 25011 Guasave & 23 \\
\midrule \textbf{Suma} & & & & \textbf{102} \\
\bottomrule
\end{tabular}
\\[2pt]\begin{minipage}{0.92\textwidth}\scriptsize Claves y AGEB, del archivo de variables; nombres, del anexo L, que marca a los siete en su columna de \emph{municipios de reciente creación}. Las 102 AGEB de este cuadro más 3 que pasan de Chinameca a Oteapan ---Veracruz, municipios que ya existían--- son las 105 que cambian de municipio; las 15 restantes de las 120 reasignadas cambian de localidad dentro del mismo municipio, absorbidas por la cabecera.\end{minipage}
\end{table}
